\subsection{Problem Statement}
We study real-time multimodal trajectory generation for robotic manipulation. Given a high-dimensional observation $\mathbf{o}_t$ consisting of multi-view RGB images, depth maps, point clouds, and proprioceptive states, the goal is to generate a set of diverse, temporally coherent future action trajectories $\{\boldsymbol{\tau}_k\}_{k=1}^K$ over a fixed horizon of $H$ steps.

The problem poses three key challenges: 
(1) Multimodal perception, requiring the robust fusion of heterogeneous sensory inputs with complementary semantic and geometric information; 
(2) Multi-hypothesis generation, where multiple plausible futures must be modeled rather than a single deterministic trajectory, effectively capturing the inherent uncertainty of manipulation tasks; and 
(3) Real-time inference, demanding high-bandwidth closed-loop control (e.g., $>100$ Hz) that fundamentally precludes iterative sampling or numerical integration procedures.

Existing diffusion- and flow-based generative policies can model complex multimodal distributions but rely heavily on iterative denoising or ODE integration (typically restricted to low frequencies like $3-10$ Hz), making them unsuitable for dynamic, real-time deployment. We aim to design a policy that preserves the multimodal expressiveness of flow models while enabling instantaneous one-step inference.

\subsection{Method Overview}
To address these challenges, we decompose real-time multimodal trajectory generation into three core components:
\begin{enumerate}
    \item A multimodal perception encoder that fuses RGB, depth, and point-cloud observations into a unified, geometry-aware state representation;
    \item A Conditional Flow Matching teacher, which acts as an expert to model a $K$-shot, $H$-step multimodal trajectory distribution in the data space using a temporal 1D U-Net;
    \item An IMLE-distilled one-step student policy, which directly approximates the teacher's distribution via a single forward pass. It adopts an identical network architecture to the teacher to strictly isolate the performance gains of the distillation algorithm.
\end{enumerate}

While the teacher captures rich multimodal behaviors, its reliance on ODE-based sampling makes it unsuitable for high-frequency control. We therefore distill the teacher into a fast-reactive student using Implicit Maximum Likelihood Estimation (IMLE). This framework requires only offline teacher samples and completely avoids score matching or multi-time supervision. A bi-directional Chamfer diversity loss further ensures that the student rigorously preserves the multi-hypothesis structure of the teacher while collapsing the entire inference process to a single step.

\begin{figure*}[t]
    \centering
    \includegraphics[width=\linewidth]{figures/Frame.pdf}
    \caption{
    Overview of our framework. We encode multimodal observations into a geometry-aware representation, train a Conditional Flow Matching teacher to model multi-hypothesis trajectories, and distill it into a one-step student policy via IMLE for real-time control.
    }
    \label{fig:framework}
\end{figure*}

\subsection{Multimodal Observation Encoding}
\label{sec:multimodal_encoding}
Our policy is conditioned on rich multimodal observations to jointly capture semantic context and explicit geometric structure. At each execution time step $t$, the observation $\mathbf{o}_t = \{\mathcal{I}_{rgb}, \mathcal{I}_{depth}, \mathcal{P}, \mathcal{S}\}$ consists of multi-view RGB images, aligned depth maps, a 3D point cloud, and the robot's proprioceptive state.

\paragraph{Visual Feature Extraction.}
RGB and depth inputs are processed by a dual-stream visual encoder. For RGB images, we employ a ResNet-18 backbone to extract dense semantic features, while depth maps are encoded using a modified ResNet-18 with a single-channel input layer to emphasize explicit geometric cues. Features extracted from multiple camera views are flattened and projected into a shared latent space, yielding the modality-specific embeddings $\mathbf{f}_{rgb}, \mathbf{f}_{depth} \in \mathbb{R}^{B \times D}$.

\paragraph{Bi-Directional Cross-Modal Attention.}
To establish dense correspondences between semantic textures and geometric structures, we introduce a bi-directional cross-modal attention mechanism. Instead of relying on naive concatenation, which struggles to capture inter-modality relationships, each modality explicitly attends to the other:
\begin{align}
    \mathbf{f}'_{rgb} &= \mathbf{f}_{rgb} + \text{Attn}(\mathbf{f}_{rgb}, \mathbf{f}_{depth}, \mathbf{f}_{depth}), \\
    \mathbf{f}'_{depth} &= \mathbf{f}_{depth} + \text{Attn}(\mathbf{f}_{depth}, \mathbf{f}_{rgb}, \mathbf{f}_{rgb}),
\end{align}
where $\text{Attn}(Q, K, V) = \text{Softmax}\left(\frac{QK^\top}{\sqrt{d}}\right)V$ denotes the scaled dot-product attention. This symmetrical formulation enriches both streams, yielding geometry-aware RGB features and semantically enhanced depth features.

\paragraph{Gated Adaptive Fusion \& Global Embedding.}
Real-world robotic environments often exhibit dynamic sensory degradation, such as lighting variations affecting RGB cameras or reflective surfaces corrupting depth sensors. To improve robustness against such modality-specific noise, a gating network adaptively predicts modality importance weights $\boldsymbol{\alpha} = [\alpha_{rgb}, \alpha_{depth}] = \text{Softmax}\bigl(G([\mathbf{f}'_{rgb}, \mathbf{f}'_{depth}])\bigr)$. The fused visual representation $\mathbf{h}_{vis}$ is then computed via the adaptively weighted sum of the projected features.

In parallel, the 3D point cloud $\mathcal{P}$ is encoded using a PointNet architecture to extract a global geometric descriptor $\mathbf{h}_{pcd}$, and the robot proprioceptive state $\mathcal{S}$ is embedded into a state vector $\mathbf{h}_{state}$ via a multi-layer perceptron. Finally, all extracted features are concatenated to form the global observation embedding:
\begin{equation}
    \mathbf{E}_{obs} = [\mathbf{h}_{vis}, \mathbf{h}_{pcd}, \mathbf{h}_{state}].
\end{equation}
This highly informative, noise-resilient representation serves as the universal condition for both the generative teacher and the distilled student policies.

\subsection{Policy Learning via Conditional Flow Matching}
Given the multimodal observation embedding $\mathbf{E}_{obs}$ (e.g., comprising RGB, depth, point cloud, and proprioception), our goal is to learn a generative policy that models a multimodal distribution over future action trajectories. To achieve this, we formulate the action generation as a Conditional Flow Matching (CFM) problem, which provides stable training dynamics and strong generative capacity in continuous domains. We define a fixed action horizon $H$ for each prediction step, effectively framing the execution as a receding-horizon control problem.

\paragraph{Flow Matching Formulation and Data-Space Objective.} 
Let $\boldsymbol{\tau}_1 \in \mathbb{R}^{H \times D_{\text{action}}}$ denote a ground-truth action trajectory sequence sampled from the expert distribution $q(\boldsymbol{\tau})$, and let $\boldsymbol{\tau}_0 \sim \mathcal{N}(\mathbf{0}, \mathbf{I})$ be a pure noise sequence of the same dimension. CFM constructs a probability path that linearly interpolates between the noise and the data:
\begin{equation}
    \boldsymbol{\tau}_t = (1 - t)\boldsymbol{\tau}_0 + t\boldsymbol{\tau}_1, \quad t \in [0,1].
\end{equation}
To train the network, we minimize the following data-space objective:
\begin{equation}
    \mathcal{L}_{\text{CFM}} = \mathbb{E}_{t, \boldsymbol{\tau}_0, \boldsymbol{\tau}_1} \left\| D_\theta(\boldsymbol{\tau}_t, t, \mathbf{E}_{obs}) - \boldsymbol{\tau}_1 \right\|^2_2.
\end{equation}

\paragraph{Logit-Normal Time Schedule and Anti-Overfitting Masking.} 
Instead of sampling the flow time $t$ uniformly, we employ a logit-normal distribution $\text{logit}(t) \sim \mathcal{N}(\mu_t, \sigma_t^2)$. This schedule concentrates more training capacity on the noisier, early-generation stages (i.e., $t \to 0$), which empirically improves the final generation quality. Furthermore, as $t \to 1$, the input $\boldsymbol{\tau}_t$ becomes highly identical to the target $\boldsymbol{\tau}_1$. To prevent the network from learning a trivial identity mapping that ignores the condition $\mathbf{E}_{obs}$, we introduce a flow-time-dependent masking mechanism. The noise embedding is stochastically masked to zero with a probability $f_m(t) = \frac{1}{1+e^{-k(t-m)}}$, acting as an embedding-level dropout that forces the model to heavily rely on the observation $\mathbf{E}_{obs}$ during the critical final stages of denoising.

\subsection{Teacher Multi-Hypothesis Trajectory Sampling}
Although the CFM teacher accurately models multimodal trajectory distributions, its inference requires iterative ODE solving, which introduces unacceptable latency for real-time high-frequency robot control. To enable distillation, we utilize the trained teacher to construct an offline synthetic dataset of multimodal futures. For each observation $\mathbf{E}_{obs}$, we sample $K$ independent noise vectors $\{\boldsymbol{\tau}_0^{(i)}\}_{i=1}^K \sim \mathcal{N}(\mathbf{0}, \mathbf{I})$ and solve the ODE using the reparameterized velocity field from $D_\theta$, yielding a set of explicit teacher predictions $\mathcal{T}_{\text{teacher}} = \{\boldsymbol{\tau}_1^{*(1)}, \dots, \boldsymbol{\tau}_1^{*(K)}\}$. This converts the implicit continuous flow density into a discrete set of high-quality trajectory modes.

\subsection{One-Step Student Policy Learning}
To break the computational bottleneck, we propose a sample-only, likelihood-free distillation framework to compress the teacher into a single-step student policy $\pi_{\text{student}}$. To rigorously isolate the performance gains of our distillation algorithm from architectural capacities, the student policy adopts an \textbf{identical 1D U-Net architecture} to the teacher. The critical modification is the complete removal of all time-step conditioning modules (e.g., the sinusoidal encoders and corresponding FiLM time-projections). Consequently, the student $\tau_\psi$ maps pure noise $\mathbf{z} \sim \mathcal{N}(\mathbf{0}, \mathbf{I})$ and the condition $\mathbf{E}_{obs}$ directly to a complete trajectory sequence in a single forward pass:
\begin{equation}
    \hat{\boldsymbol{\tau}} = \tau_\psi(\mathbf{E}_{obs}, \mathbf{z}) \in \mathbb{R}^{H \times D_{\text{action}}}.
\end{equation}
During closed-loop deployment, the student generates the fixed $H$-step horizon at $>100$ Hz, but only the first $T_e$ steps (Execution Horizon) are executed on the robot before re-planning with updated observations, thereby forming a highly reactive receding-horizon control loop.

\subsection{Set-Level IMLE Distillation via Chamfer Matching}
Standard distillation objectives (e.g., MSE or KL-divergence) often lead to mode collapse, where the student merely averages the multimodal outputs of the teacher into a single, sometimes physically implausible, mean trajectory. To address this, we formulate the distillation through the lens of Implicit Maximum Likelihood Estimation (IMLE), which inherently guarantees mode preservation.

During training, the student model generates a set of $K$ candidate hypotheses $\{\hat{\boldsymbol{\tau}}_j\}_{j=1}^K$ conditioned on the current observation. To align this generated distribution directly with the $K$ explicit modes provided by the teacher's trajectory set $\mathcal{T}_{\text{teacher}} = \{\boldsymbol{\tau}_i^*\}_{i=1}^K$, we employ a bi-directional Chamfer distance as a set-level proxy for the IMLE objective:
\begin{align}
    \mathcal{L}_{\text{Chamfer}} &= \frac{1}{K} \sum_{i=1}^K \min_{j \in [K]} \left\| \boldsymbol{\tau}_i^* - \hat{\boldsymbol{\tau}}_j \right\|^2_2 \nonumber \\
    &\quad + \frac{1}{K} \sum_{j=1}^K \min_{i \in [K]} \left\| \boldsymbol{\tau}_i^* - \hat{\boldsymbol{\tau}}_j \right\|^2_2 .
\end{align}
This symmetric formulation serves two critical purposes:
\begin{itemize}
    \item \textbf{Mode-Covering:} The first term ensures that \textit{every} multimodal behavior predicted by the teacher is successfully covered by at least one trajectory from the student's candidate pool, thereby strictly preventing mode collapse.
    \item \textbf{Mode-Seeking:} The second term acts as a structural regularizer, ensuring that all trajectories generated by the student are grounded in the valid, physically plausible behavior manifold defined by the expert teacher.
\end{itemize}
By minimizing this set-level matching objective, the single-step student policy faithfully reconstructs the full geometric and statistical diversity of the teacher's flow distribution, requiring only a fraction of the computational cost.